% 1.引言.tex —— 一、引言
\section{引言}

\subsection{问题背景}
干燥是影响中药材成品品质的关键工序。热风烘干通过调控烘房温湿环境完成药材升温与脱水，内部温度和含水率分布决定了局部干燥进度。仅依据表面状态或整体平均含水率，可能掩盖中心尚未达标的情况，因此需要同时分析径向差异与全场达标时间。

预热期温度逐步接近环境，水分则从内部向表面迁移。随着干燥持续，热物性和水分扩散能力发生变化，短时间常热物性模型需要扩展为变物性模型。本文以题设参数为依据，利用数值模拟连接预热分析、全过程演化和终止时间判断，并区分模型假设、数值误差与输出精度。

\subsection{问题重述}
某中药材近似为长 $25$ cm、半径 $2$ cm 的圆柱，初始温度为 $28$ $^\circ$C、干基含水率为 $2.55$ kg/kg，烘房环境记录见附件 1。四个问题的要求如下。

\textbf{问题一：}依据附录 2 建立预热平衡阶段模型，按表 1、表 2 给出 $100,300,600,900,1200,1500,1800$ s、距中心 $0,0.5,1,1.5,2$ cm 处的温度与水分浓度，并将 $1800$ s 内每隔 $1$ s、每隔 $0.1$ cm 的完整结果保存到 \texttt{result1.xlsx}，结果保留四位小数。

\textbf{问题二：}依据附录 3 的变物性经验关系，建立持续 $2\sim3$ 天的全烘干过程模型。按表 3、表 4 给出前 $3$ h 内每隔 $0.5$ h、上述五个径向位置的温度和水分浓度，并将完整结果保存到 \texttt{result2.xlsx}。本文将全过程模拟至 $72$ h，前 $3$ h 的表格用于展示预热期响应。

\textbf{问题三：}沿用问题二的参数，确定药材各处水分浓度均低于 $0.15$ kg/kg 所需的烘干时间。按表 5 给出每隔 $6$ h、径向每隔 $0.5$ cm 的水分浓度，并将每隔 $60$ s、径向每隔 $0.1$ cm 的完整结果保存到 \texttt{result3.xlsx}。终止判据取全场最大值，同时考察时间推进与空间网格对达标时间的影响。

\textbf{问题四：}考虑烘干过程中因水分流失导致的尺寸变化（收缩），依据附录 4 的经验公式与附件 2 的收缩半径数据，重新确定烘干时长。按表 6 给出每隔 $6$ h、到中心距离 $0,0.5,1.0,1.5$ cm 与药材表面处的水分浓度，并将每隔 $60$ s、径向每隔 $0.1$ cm 的完整结果保存到 \texttt{result4.xlsx}。与前三个问题的差别在于半径不再固定，扩散区域成为随时间移动的边界。
